\section{Introdução}
\label{sec:introducao}

\subsection{Contexto}

A Internet das Coisas (IoT) liga sensores, atuadores e serviços de
\emph{back-end} que trocam pequenas mensagens, com frequência e, muitas vezes, a
partir de dispositivos com recursos muito limitados. Entre os protocolos da camada
de aplicação usados neste domínio~\cite{dizdarevic2019survey}, o
\emph{Message Queuing Telemetry Transport} (MQTT) impôs-se pela sua simplicidade:
normalizado pela OASIS nas versões 3.1.1~\cite{oasis-mqtt311} e
5.0~\cite{oasis-mqtt5}, segue o paradigma publicação/subscrição, no qual os
clientes nunca comunicam diretamente uns com os outros. Um \emph{publisher}
envia mensagens para um \emph{tópico} (por exemplo,
\texttt{fabrica/sala1/temperatura}), e um intermediário, o \emph{broker}, entrega-as
a todos os clientes cuja subscrição corresponda a esse tópico. Implementações
abertas como o Eclipse Mosquitto~\cite{light2017mosquitto} tornaram o protocolo
omnipresente, de casas inteligentes a ambientes industriais e médicos.

\subsection{Problema}

A especificação do MQTT trata a segurança como responsabilidade da
implementação: oferece um campo de nome de utilizador e palavra-passe no pedido de
ligação e recomenda o uso de TLS, mas não define proteção da mensagem em si.
Daqui resultam três problemas que estruturam este trabalho.

\begin{enumerate}[label=\textbf{P\arabic*.}, leftmargin=*]
  \item \textbf{O \emph{broker} é um ponto de confiança total.} O TLS protege
    apenas cada salto cliente--\emph{broker}~\cite{rfc8446}. No \emph{broker}, as
    mensagens existem em claro: quem o comprometer, ou qualquer cliente a quem
    uma lista de controlo de acessos mal configurada permita subscrever o filtro
    \texttt{\#}, vê e pode reencaminhar todo o tráfego~\cite{andy2017attack,
    segarra2020mqttz}.
  \item \textbf{Uma mensagem autêntica pode ser repetida ou desviada.} Mesmo uma
    mensagem criptograficamente autêntica pode ser capturada e repetida mais tarde,
    ou publicada num tópico diferente daquele para que foi criada. A autenticação
    não impede estes ataques: são precisos, em separado, um carimbo temporal
    autenticado (validade temporal), números de sequência (anti-repetição) e a
    ligação criptográfica ao tópico.
  \item \textbf{Comprometimento de dispositivos.} Num parque com milhares de
    dispositivos, alguns acabarão comprometidos. Se uma chave válida puder
    publicar em qualquer tópico, um único dispositivo comprometido pode fazer-se
    passar por todos os outros, e o incidente de botnets como a
    Mirai~\cite{kolias2017mirai} mostra a escala que estes ataques podem atingir.
\end{enumerate}

A literatura propõe numerosas soluções parciais, desde esquemas de autenticação
e autorização a cifras baseadas em atributos e ambientes de execução confiáveis,
mas é difícil perceber que propriedades cada uma garante, que ameaças deixa por
tratar e como se combinam. É este o problema que a revisão endereça:
\emph{que mecanismos de segurança propõe o estado da arte para comunicações
baseadas em MQTT, que propriedades garantem e que lacunas persistem?}

\subsection{Objetivos}

\begin{enumerate}[label=\textbf{O\arabic*.}, leftmargin=*]
  \item Conduzir uma revisão sistemática da literatura, segundo um protocolo
    explícito e reprodutível~\cite{kitchenham2007, page2021prisma}, sobre
    mecanismos de segurança para o MQTT.
  \item Classificar as soluções encontradas numa taxonomia e compará-las segundo
    um conjunto comum de propriedades de segurança e de avaliação.
  \item Comparar a presente revisão com as revisões já publicadas sobre o tema.
  \item Identificar as lacunas do estado da arte e, a partir delas, propor uma
    arquitetura de segurança em camadas, descrita formalmente.
  \item Definir o objetivo e o plano experimental da segunda parte do projeto:
    medir, num dispositivo restrito (Raspberry Pi), o custo de desempenho e o
    benefício de segurança de cada camada, para identificar a configuração com a
    maior proteção e o menor impacto no desempenho.
\end{enumerate}

\subsection{Contributos}

\begin{itemize}[leftmargin=*]
  \item Um protocolo de revisão completo (perguntas, PICOC, \emph{strings} por base
    de dados, critérios de inclusão, exclusão e qualidade, formulário de
    extração), reutilizável e reprodutível (Secção~\ref{sec:metodologia}).
  \item Uma comparação de \nIncluidos{} estudos primários e de \nSecundarios{}
    revisões relacionadas, com uma matriz de propriedades de segurança
    (Secções~\ref{sec:relacionados} e~\ref{sec:resultados}).
  \item A formalização de mecanismos que a literatura raramente combina: ligação
    da mensagem ao tópico, janela anti-repetição, política do recetor contra a
    redução do nível de proteção e restrição chave--tópico
    (Secção~\ref{sec:proposta}).
  \item Uma arquitetura modular cuja implementação de referência é disponibilizada
    num repositório público~\cite{mqttsecure2026repo}, com documentação técnica
    completa.
\end{itemize}

\subsection{Estrutura do documento}

A Secção~\ref{sec:enquadramento} apresenta os fundamentos: o protocolo MQTT, os
seus mecanismos nativos de segurança, as primitivas criptográficas relevantes e o
modelo de ameaças. A Secção~\ref{sec:metodologia} descreve o protocolo da revisão
sistemática. A Secção~\ref{sec:relacionados} analisa as revisões já publicadas e
justifica a presente. A Secção~\ref{sec:resultados} apresenta os resultados: a
caracterização estatística dos estudos, a taxonomia, a comparação detalhada e as
respostas às perguntas de investigação. A Secção~\ref{sec:proposta} propõe a
arquitetura em camadas e posiciona-a face aos trabalhos relacionados. A
Secção~\ref{sec:plano} define o objetivo e o plano experimental da segunda parte
do projeto. A Secção~\ref{sec:discussao} discute as lacunas e as ameaças à validade, e a
Secção~\ref{sec:conclusoes} conclui e aponta o trabalho futuro. Os anexos contêm as
\emph{strings} de pesquisa adaptadas a cada base de dados e a extração de dados por
estudo.
